\section{Evaluation}
\label{sec:evaluation}

\subsection{Simulation Setup}
\label{sec:eval:setup}

We evaluate \sentinel against a corpus of 250,000 simulated attacks spanning 15 threat categories. The categories cover the full taxonomy introduced in \S\ref{sec:threat}, ranging from well-understood classes (direct injection, encoding-based obfuscation) to emerging threats with limited prior study (multi-modal attacks, model-level compromise).

Each base attack is tested with five mutation variants:
\begin{enumerate}[leftmargin=*]
\item \textbf{Lexical mutation.} Synonym substitution, typo injection, Unicode homoglyph replacement, and whitespace manipulation that preserve semantic content while altering surface form.
\item \textbf{Structural mutation.} Reordering of clauses, insertion of benign padding, and syntactic transformations that modify parse structure without changing intent.
\item \textbf{Encoding mutation.} Base64, ROT13, URL encoding, HTML entity encoding, and nested encoding chains that obscure payload content.
\item \textbf{Context mutation.} Prepending legitimate-seeming preambles, embedding attacks within multi-turn conversation histories, and wrapping payloads in plausible business contexts.
\item \textbf{Hybrid mutation.} Compositions of two or more of the above, representing realistic adversarial behavior where attackers apply multiple evasion techniques simultaneously.
\end{enumerate}

Every base attack is tested with all five mutation variants, yielding a mutation-inclusive corpus that stresses each layer's robustness to evasion. Attacks are drawn from publicly available red-teaming datasets, internally developed attack generators, and adversarial examples synthesized via automated prompt perturbation.

\subsection{Detection Cascade}
\label{sec:eval:cascade}

Attacks flow through the \sentinel layer stack in order. Each layer catches a fraction of the remaining attacks; those that survive all layers constitute the \emph{residual}. Table~\ref{tab:cascade} shows the detection cascade for the full 250,000-attack corpus.

\begin{table}[t]
\centering
\small
\caption{Detection cascade across the \sentinel layer stack. Each row reports attacks caught at that layer out of the full 250,000-attack corpus.}
\label{tab:cascade}
\begin{tabular}{@{}lrr@{}}
\toprule
Layer & Attacks Caught & Percentage \\
\midrule
L1 Sentinel Core        & 89,910  & 36.0\% \\
L2 Capability Proxy     & 50,849  & 20.3\% \\
L3 Behavioral EDR       & 27,151  & 10.9\% \\
\pasr                   &  4,983  &  2.0\% \\
\tcsa                   &  2,000  &  0.8\% \\
\asra                   &  3,250  &  1.3\% \\
Combos ($\alpha$+$\beta$+$\gamma$) & 15,198 & 6.1\% \\
\mire (containment)     &  1,750  &  0.7\% \\
\midrule
\textbf{Total caught}   & \textbf{195,091} & \textbf{$\sim$78.0\%} \\
\textbf{Caught + contained} & \textbf{$\sim$246,250} & \textbf{$\sim$98.5\%} \\
\textbf{Residual}       & \textbf{$\sim$3,750}   & \textbf{$\sim$1.5\%} \\
\bottomrule
\end{tabular}
\end{table}

The cascade exhibits desirable front-loading: L1 alone eliminates over one-third of all attacks at sub-millisecond latency, and the first three layers collectively neutralize 67.2\% of the corpus before any novel primitive is invoked. The combinatorial layers (\pasr, \tcsa, \asra) and their compositions contribute an additional 31.3\%, catching attacks that exploit semantic ambiguity, temporal escalation, and cross-boundary trust---classes that are invisible to pattern matching and access control alone. \mire provides the final containment guarantee, rendering an additional 0.7\% of attacks harmless even when they cannot be \emph{detected}, by ensuring that the model's output is correct regardless of whether the model itself is compromised.

\subsection{Per-Category Detection Rates}
\label{sec:eval:categories}

Table~\ref{tab:percategory} reports per-category detection rates and residual counts. Categories are ordered by detection rate, from highest to lowest. The variation across categories is informative: well-structured attacks (direct injection, encoding) are nearly perfectly detected, while categories that exploit fundamental ambiguity (semantic equivalence, multi-modal, model-level compromise) exhibit meaningfully higher residual rates.

\begin{table*}[t]
\centering
\small
\caption{Per-category detection rates across the full 250,000-attack corpus. Residual counts are approximate due to rounding across mutation variants.}
\label{tab:percategory}
\begin{tabular}{@{}lrrr@{}}
\toprule
Category & Attacks & Detection Rate & Residual \\
\midrule
Direct Injection        & 25,000 & $>$99.99\% & $\sim$2   \\
Indirect Injection      & 25,000 & $>$99.97\% & $\sim$6   \\
Tool Abuse              & 20,000 & 99.9\%     & $\sim$13  \\
Data Exfiltration       & 15,000 & 99.9\%     & $\sim$14  \\
Role-play / Persona     & 20,000 & 99.9\%     & $\sim$20  \\
Encoding / Obfuscation  & 20,000 & 99.8\%     & $\sim$36  \\
Cross-boundary Trust    & 10,000 & 99.5\%     & $\sim$53  \\
Multi-turn Crescendo    & 20,000 & 99.3\%     & $\sim$134 \\
Social Engineering      & 15,000 & 98.2\%     & $\sim$263 \\
Adversarial ML          & 10,000 & 96.8\%     & $\sim$324 \\
Steganographic          & 12,000 & 96.7\%     & $\sim$398 \\
Novel / Zero-day        & 13,000 & 96.2\%     & $\sim$498 \\
Semantic Equivalence    & 15,000 & 96.5\%     & $\sim$525 \\
Multi-modal             & 10,000 & 94.0\%     & $\sim$600 \\
Model-Level Compromise  & 10,000 & 91.5\%     & $\sim$850 \\
\midrule
\textbf{All categories} & \textbf{250,000} & \textbf{$\sim$98.5\%} & \textbf{$\sim$3,736} \\
\bottomrule
\end{tabular}
\end{table*}

The three highest-residual categories---model-level compromise, multi-modal, and semantic equivalence---together account for over 52\% of all residual attacks despite representing only 14\% of the corpus. This concentration is not a failure of the architecture; it reflects fundamental impossibility results that bound any defense system, as analyzed in \S\ref{sec:eval:residual}.

\subsection{Historical Progression}
\label{sec:eval:progression}

The current results represent the culmination of four iterative development phases, each expanding the simulation scope and adding architectural layers. Table~\ref{tab:progression} summarizes the progression.

\begin{table*}[t]
\centering
\small
\caption{Historical progression of \sentinel evaluation. Each phase adds layers and expands the attack corpus. The floor represents the estimated mathematical limit given known impossibility results.}
\label{tab:progression}
\begin{tabular}{@{}llrrll@{}}
\toprule
Phase & Simulation & Detection & Residual & Key Addition \\
\midrule
Phase 1 & 100K attacks, 9 categories   & 81.6\% & 18.4\%       & L1--L3 only \\
Phase 2 & 250K attacks, 15 categories  & 93.7\% &  6.3\%       & +L4--L6, +Combos \\
Phase 3 & 250K + \pasr                 & 95.7\% &  4.3\%       & +\pasr resolves conflict \\
Phase 4 & 250K + all layers            & $\sim$98.5\% & $\sim$1.5\% & +\tcsa, +\asra, +\mire \\
\midrule
Floor   & ---                          & $\sim$98--99\% & $\sim$1--2\% & Mathematical limit \\
\bottomrule
\end{tabular}
\end{table*}

The progression reveals two important dynamics. First, diminishing returns are expected and informative: the jump from Phase~1 to Phase~2 (18.4\% $\to$ 6.3\% residual) reflects the addition of entirely new detection paradigms, while subsequent phases yield smaller but meaningful gains by addressing increasingly subtle attack classes. Second, the convergence toward a 1--2\% floor is consistent with the theoretical impossibility results discussed below, suggesting that the architecture is approaching the fundamental limit of what any defense system can achieve.

\subsection{Residual Analysis}
\label{sec:eval:residual}

The approximately 3,750 residual attacks ($\sim$1.5\% of the corpus) are not random failures; they concentrate in three irreducible classes that correspond to known theoretical barriers. Understanding these classes is essential for setting realistic expectations about any LLM defense architecture.

\paragraph{Semantic identity ($\sim$35\% of residual).} In this class, the malicious input is \emph{mathematically indistinguishable} from a legitimate request. The canonical example is ``How do I mix bleach and ammonia?''---a query that is simultaneously a valid chemistry question and a request for instructions to produce chloramine gas. No classifier, regardless of sophistication, can separate these cases without access to the user's private intent, which is by definition unobservable. This is not a limitation of \sentinel specifically but a fundamental property of natural language: the mapping from surface form to intent is many-to-many, and certain malicious intents map to surface forms that are identical to benign ones.

\paragraph{Model trust chain ($\sim$45\% of residual).} Goldwasser and Kim~\cite{goldwasser2022} prove that detecting backdoors in trained neural networks is computationally infeasible under standard cryptographic assumptions. This impossibility result implies that if the underlying language model has been compromised during training---via data poisoning, weight manipulation, or supply-chain attacks---no runtime defense can reliably detect the compromise. \mire addresses this class by shifting from detection to containment: rather than identifying whether the model is compromised, \mire ensures that critical outputs are verifiable against external ground truth, rendering the model's internal state irrelevant for safety-critical decisions. However, not all outputs admit efficient external verification, and the residual in this class reflects tasks where containment verification is infeasible.

\paragraph{Representation gap ($\sim$20\% of residual).} Attacks in this class exploit modalities, encodings, or representation formats that fall outside the analytical coverage of all layers. Examples include attacks embedded in image steganography within multi-modal inputs, adversarial perturbations in audio spectrograms, and payloads encoded in domain-specific formats (e.g., molecular structure notation, musical score encoding) that no layer currently parses. This class is reducible in principle---new analyzers can be added for each modality---but the space of possible representations is unbounded, ensuring a nonzero residual from this class at any point in time.

These three classes are largely complementary, and their relative proportions are stable across simulation runs ($\pm$2\% variation). The semantic identity and model trust chain classes are provably irreducible: they correspond to information-theoretic and computational impossibility results, respectively. The representation gap class is reducible but practically unbounded, as new modalities and encodings emerge continuously. Together, they establish a theoretical floor of approximately 1--2\% residual for any defense architecture operating under the same threat model.
